\subsection{COMPACT SETS OF CONTINUOUS MAPPINGS}

THEOREM 2 (Ascoli). Let X be a topological (resp. uniform) space, let S be a covering of X, let Y be a uniform space and H a set of mappings of X into Y such that, for each A ∈ S and each u ∈ H, the restriction of u to A is continuous (resp. uniformly continuous). Then, for H to be precompact with respect to the uniformity of S-convergence, it is necessary in all cases and also sufficient if the sets \(\mathbf{A} \in \mathfrak{S}\) are compact (resp. precompact) that the following conditions should be satisfied:

\begin{enumerate}
\def\labelenumi{\alph{enumi})}
\item
  For each \(\mathbf{A} \in \mathfrak{S}\) , the set \(\mathbf{H}|\mathbf{A} \subset \mathcal{F}(\mathbf{A};\mathbf{Y})\) of restrictions to \(\mathbf{A}\) of functions of \(\mathbf{H}\) is equicontinuous (resp. uniformly equicontinuous).
\item
  For each \(x \in \mathbf{X}\) , the set \(\mathbf{H}(x) \subset \mathbf{Y}\) of points \(u(x)\) ( \(u \in \mathbf{H}\) ) is precompact.
\end{enumerate}

\begin{enumerate}
\def\labelenumi{\arabic{enumi})}
\tightlist
\item
  Let us show first that conditions a) and b) are necessary. We know (§ 1, no. 2, Remark 6) that the mapping \(u \to u(x)\) of \(\mathcal{F}_{\mathfrak{S}}(\mathbf{X};\mathbf{Y})\) into Y is uniformly continuous; hence, if H is precompact, so is \(\mathbf{H}(x)\) (Chapter II, § 4, no. 2, Proposition 2), which proves b). To prove a), consider a set \(\mathbf{A} \in \mathfrak{S}\) , a point \(x_0 \in \mathbf{A}\) and an entourage V of Y; since H is precompact it can be covered by a finite number of \(W(\mathbf{A},\mathbf{V})\) -small sets; in other words there is a finite sequence \((u_i)\) of elements of H such that, for each \(u \in \mathbf{H}\) , we have
\end{enumerate}

\[
(u (x), u _ {i} (x)) \in \mathbf {V} \quad \text { for   all } \quad x \in \mathbf {A}\tag{4}
\]

for at least one index i.

Since each of the \(u_{i}|A\) is continuous at \(x_{0}\) (resp. uniformly continuous) there is a neighbourhood \(U_{i}\) of \(x_{0}\) in A (resp. an entourage \(M_{i}\) of A) such that

\[
x \in \mathrm{U} _ {i} \quad \text { implies } \quad (u _ {i} (x), u _ {i} (x _ {0})) \in \mathrm{V},\tag{5}
\]

(resp. such that

\[
(x ^ {\prime}, x ^ {\prime \prime}) \in \mathbf {M} _ {i} \quad \text { implies } \quad (u _ {i} (x ^ {\prime}), u _ {i} (x ^ {\prime \prime})) \in \mathbf {V}.\tag{6}
\]

Let U (resp. M) be the intersection of the \(U_{i}\) (resp. \(M_{i}\) ); it is a neighbourhood of \(x_{0}\) in A (resp. an entourage of A). For each \(u \in H\) there is an index i for which (4) holds; writing condition (4) for \(x_{0}\) and for x (resp. for \(x'\) and \(x''\) ) and taking account of (5) {[}resp. (6){]}, we see immediately that the relation \(x \in U\) {[}resp. \((x', x'') \in M\) {]} implies \((u(x), u(x_{0})) \in \overset{\mathbf{3}}{\mathrm{V}}\) {[}resp. \((u(x'), u(x'')) \in \overset{\mathbf{3}}{\mathrm{V}}\) {]}, for each \(u \in H\) ; and this establishes a).

\begin{enumerate}
\def\labelenumi{\arabic{enumi})}
\setcounter{enumi}{1}
\tightlist
\item
  Now let us show that the conditions a) and b) are sufficient if the sets \(A \in S\) are compact (resp. precompact). Condition b) implies that H is precompact with respect to the uniformity of pointwise convergence (Chapter II, § 4, no. 2, Proposition 3). But it follows from condition a) and Theorem 1 of no. 4 that on H\textbar A the uniformity of pointwise convergence in A coincides with the uniformity of uniform convergence in A; hence H\textbar A is precompact in \(\mathcal{F}_{u}(A; Y)\) , which implies that H is precompact with respect to the uniformity of S-convergence (§ 1, no. 2).
\end{enumerate}

Note that condition b) of Theorem 2 is automatically satisfied if Y is a precompact space.

COROLLARY 1. Let X be a topological (resp. uniform) space, let Y be a Hausdorff uniform space and let H be an equicontinuous (resp. uniformly equicontinuous) subset of \(\mathcal{C}(\mathbf{X};\mathbf{Y})\) . Suppose that \(\mathrm{H}(x)\) is relatively compact in Y for each \(x\in \mathbf{X}\) . Then H is relatively compact in \(\mathcal{C}(\mathbf{X};\mathbf{Y})\) with respect to the topology of compact (resp. precompact) convergence.

Let \(\overline{\mathbf{H}}\) be the closure of \(\mathbf{H}\) in \(\mathcal{F}_s(\mathbf{X};\mathbf{Y})\) . \(\overline{\mathbf{H}}\) is equicontinuous (resp. uniformly equicontinuous) (no. 3, Proposition 6). Moreover, we have \(\overline{\mathbf{H}}(x)\subset \overline{\mathbf{H}(x)}\) ( \(\S\) 1, no. 2, Remark 6) and therefore \(\overline{\mathbf{H}}(x)\) is also relatively compact; hence Theorem 2 shows that \(\overline{\mathbf{H}}\) is precompact with respect to \(\mathfrak{S}\) -convergence, where \(\mathfrak{S}\) denotes the set of all compact (resp. precompact) subsets of \(\mathbf{X}\) . Moreover, since \(\overline{\mathbf{H}(x)}\) is compact, and therefore complete, \(\overline{\mathbf{H}}\) is complete with respect to the uniformity of pointwise convergence (Chapter II, \(\S\) 3, no. 5, Proposition 10 and no. 4, Proposition 8) and therefore also with respect to the uniformity of \(\mathfrak{S}\) -convergence ( \(\S\) 1, no. 5, Proposition 5, Corollary 2); \(\overline{\mathbf{H}}\) is therefore compact, since it is precompact, complete and Hausdorff ( \(\S\) 1, no. 2, Proposition 1).

COROLLARY 2. Let X be a topological (resp. uniform) space, let Y be a complete Hausdorff uniform space and let H be an equicontinuous (resp. uniformly equicontinuous) subset of C(X; Y). Suppose that H(x) is relatively compact in Y for all \(x \in D\) , where D is a dense subset of X. Then H is relatively compact in C(X; Y) with respect to the topology of compact (resp. precompact) convergence.

It is enough to show that \(\mathrm{H}(x)\) is relatively compact for all \(x \in X\) , for we can then apply Corollary I. Since Y is complete it is enough to show that \(\mathrm{H}(x)\) is precompact for all \(x \in X\) . Now if V is any symmetric entourage of Y, there is a neighbourhood U of x such that \((u(x), u(x')) \in V\) for all \(x' \in U\) and all \(u \in H\) . By hypothesis there exists \(x' \in U \cap D\) , and since \(\mathrm{H}(x')\) is relatively compact in Y, there exists a finite number of points \(y_k \in Y\) such that \(\mathrm{H}(x')\) is contained in the union of the sets \(\mathrm{V}(y_k)\) ; hence \(\mathrm{H}(x)\) is contained in the union of the sets \(\overset{2}{\mathrm{V}}(y_k)\) , and the proof is complete.

COROLLARY 3. Let X be a locally compact space, Y a Hausdorff uniform space, H a subset of C (X; Y). Then H is relatively compact in \(\mathcal{C}_c(X; Y)\) if and only if H is equicontinuous and H(x) relatively compact in Y for all \(x \in \mathbf{X}\) .

In view of Corollary 1 it is enough to show that, if H is relatively compact in \(\mathcal{C}_{c}(\mathbf{X};\mathbf{Y})\) , then H is equicontinuous. Now each point \(x\in X\) has a compact neighbourhood A, and it follows from Theorem 2 that H/A is equicontinuous; this implies that H is equicontinuous at x, and the result is proved.

Remark. Let X be a topological space, Y a uniform space and S a set of subsets of X. Then on every precompact subset H of \(\mathcal{F}_{\mathfrak{S}}(X;Y)\) , the uniformity of S-convergence is the same as the uniformity of pointwise convergence in \(B = \bigcup_{A \in \mathfrak{S}} A\) . We can reduce to the case where \(B = X\) and Y is Hausdorff and complete; for if \(j\) is the canonical injection \(B \to X\) and \(i\) the canonical mapping \(Y \to \hat{Y}\) , the uniformity of S-convergence on \(\mathcal{F}(X;Y)\) is the inverse image of the uniformity of S-convergence on \(\mathcal{F}(B;\hat{Y})\) under the mapping \(\theta : u \to i \circ u \circ j\) ( \(\S 1\) , no. 4, Proposition 4), and H is precompact if and only if \(\theta(H)\) is (Chapter II, \(\S 4\) , no. 2, Proposition 3). This being so, if \(B = X\) and Y is Hausdorff and complete, \(\mathcal{F}_{\mathfrak{S}}(X;Y)\) is Hausdorff and complete ( \(\S 1\) , no. 2, Proposition 1 and no. 5, Theorem 1); hence the closure \(\overline{H}\) of H in this space is compact. On \(\overline{H}\) , the topology of pointwise convergence is Hausdorff ( \(\S 1\) , no. 2, Proposition 1) and coarser than that of S-convergence; hence these two topologies coincide (Chapter I, \(\S 9\) , no. 4, Theorem 2, Corollary 3) and consequently so do the uniformities of S-convergence and pointwise convergence (Chapter II, \(\S 4\) , no. 1, Theorem 1).

